
\documentclass{article}
\usepackage{colortbl}
\usepackage{makecell}
\usepackage{multirow}
\usepackage{supertabular}

\begin{document}

\newcounter{utterance}

\centering \large Interaction Transcript for game `dond', experiment `semi\_en', episode 1 with reflexion{-}Qwen3.5{-}4B.
\vspace{24pt}

{ \footnotesize  \setcounter{utterance}{1}
\setlength{\tabcolsep}{0pt}
\begin{supertabular}{c@{$\;$}|p{.15\linewidth}@{}p{.15\linewidth}p{.15\linewidth}p{.15\linewidth}p{.15\linewidth}p{.15\linewidth}}
    \# & $\;$A & \multicolumn{4}{c}{Game Master} & $\;\:$B\\
    \hline

    \theutterance \stepcounter{utterance}  
    & & \multicolumn{4}{p{0.6\linewidth}}{
        \cellcolor[rgb]{0.9,0.9,0.9}{
            \makecell[{{p{\linewidth}}}]{
                \texttt{\tiny{[P1$\langle$GM]}}
                \texttt{You have attempted this same task before and did not succeed. These are the plans you wrote after those attempts.} \\
\texttt{Trial \#1: New plan: I failed because I made a unilateral proposal without first confirming the other player's valuation or checking if their preferences align with mine. In this game, the other player's secret value function is unknown, so I cannot assume they will accept a deal that leaves them with low{-}value items (like the eraser and battery) if they value the towels highly. My strategy should be to first communicate my valuation and ask the other player to reveal theirs (or make a counter{-}proposal) to ensure our proposals are complementary. Specifically, I should propose a split that maximizes the sum of our values, such as me taking the two towels (8 points) and the other player taking the bell and eraser (1+1=2 points), but only after verifying that the other player does not value the towels more than the bell/eraser combination. If I propose directly without this check, the other player might reject it or make a proposal that doesn't sum to a valid exchange, resulting in zero points. Next time, I will start by stating my values and asking the other player to state theirs, or propose a split that explicitly accounts for the possibility that the other player values the towels even higher, ensuring we reach a mutually acceptable agreement before finalizing the secret proposal.} \\
\\ 
\\ 
\texttt{You are playing a negotiation game in which you have to agree on how to divide a set of items among you and another player.} \\
\\ 
\texttt{Rules:} \\
\texttt{(a) You and the other player are given a set of items. Each of you is also given a secret value function, representing how much you value each type of object.} \\
\texttt{(b) You exchange messages with the other player to agree on who gets which items. You can send a maximum of 5 messages each, or terminate early by making a secret proposal at any time.} \\
\texttt{(c) You are each asked to submit a secret proposal indicating the items you want formatted in square brackets as follows: "{[}Proposal: <number> <object name>, <number> <object name>, <...>{]}"} \\
\texttt{(d) If your proposals are complementary, i.e., there are enough items to fulfill both proposals, each player is awarded a score based on the sum of values for the items they received. Otherwise, both of you get zero points.} \\
\texttt{(e) Your goal is to maximize the score you receive.} \\
\\ 
\texttt{Let us start.} \\
\\ 
\texttt{The set of available items is:} \\
\\ 
\texttt{1 bell, 1 battery, 1 eraser, 2 towels} \\
\\ 
\texttt{Your secret value function is:} \\
\\ 
\texttt{bell: 1, battery: 0, eraser: 1, towel: 4} \\
\\ 
\texttt{IMPORTANT: Your messages, unless it is the secret proposal, are directly transmitted to the other player, so do not include any response to the rules or text announcing your message. To make a secret proposal, use the indicated format. Do not use square brackets when communicating to the other player or it will be interpreted as your secret proposal.} \\
            }
        }
    }
    & \\ \\

    \theutterance \stepcounter{utterance}  
    & \multicolumn{4}{p{0.6\linewidth}}{
        \cellcolor[rgb]{0.9,0.9,0.9}{
            \makecell[{{p{\linewidth}}}]{
                \texttt{\tiny{[P1$\rangle$GM]}}
                \texttt{I value the towels very highly (4 points each), while the bell and eraser are worth 1 point each, and the battery is worthless to me. To maximize our total value, I should take the two towels (8 points) and leave the bell, eraser, and battery for the other player. However, since I don't know how much the other player values the towels, I must first communicate my valuation to ensure they don't value the towels more than the other items.} \\
\\ 
\texttt{I will state my values and ask for theirs to find a mutually beneficial split.} \\
\\ 
\texttt{My values: bell: 1, battery: 0, eraser: 1, towel: 4. What are your values?} \\
            }
        }
    }
    & & \\ \\

    \theutterance \stepcounter{utterance}  
    & & & \multicolumn{2}{p{0.3\linewidth}}{
        \cellcolor[rgb]{0.9,0.9,0.9}{
            \makecell[{{p{\linewidth}}}]{
                \texttt{\tiny{[GM$|$GM]}}
                \texttt{continue} \\
            }
        }
    }
    & & \\ \\

    \theutterance \stepcounter{utterance}  
    & & & \multicolumn{2}{p{0.3\linewidth}}{
        \cellcolor[rgb]{0.9,0.9,0.9}{
            \makecell[{{p{\linewidth}}}]{
                \texttt{\tiny{[GM$|$GM]}}
                \texttt{I value the towels very highly (4 points each), while the bell and eraser are worth 1 point each, and the battery is worthless to me. To maximize our total value, I should take the two towels (8 points) and leave the bell, eraser, and battery for the other player. However, since I don't know how much the other player values the towels, I must first communicate my valuation to ensure they don't value the towels more than the other items.} \\
\\ 
\texttt{I will state my values and ask for theirs to find a mutually beneficial split.} \\
\\ 
\texttt{My values: bell: 1, battery: 0, eraser: 1, towel: 4. What are your values?} \\
            }
        }
    }
    & & \\ \\

    \theutterance \stepcounter{utterance}  
    & & \multicolumn{4}{p{0.6\linewidth}}{
        \cellcolor[rgb]{0.9,0.9,0.9}{
            \makecell[{{p{\linewidth}}}]{
                \texttt{\tiny{[GM$\rangle$P2]}}
                \texttt{You have attempted this same task before and did not succeed. These are the plans you wrote after those attempts.} \\
\texttt{Trial \#1: New plan: I failed because I submitted a proposal for only two items (1 battery, 1 bell) while there are four items in total. Since the rules state that proposals must be complementary to award any points, my proposal left the other player with no items to offer, making a valid agreement impossible. I must ensure my proposal includes enough items to cover the other player's potential needs based on the total set. Given my values (battery: 3, bell: 2, towel: 2, eraser: 1), I should propose taking the highest{-}value items I can while leaving the remaining items for the other player to form a valid set. A better strategy is to propose taking the battery and one towel (total value 5), leaving the bell and eraser for the other player (total value 3), or taking the battery and bell (value 5) and leaving the two towels (value 4). I will calculate the total value of the items I leave behind to ensure it is at least as high as what I get, or simply ensure the remaining items form a complete set for the opponent. Specifically, I should propose: {[}Proposal: 1 battery, 1 towel{]} to leave {[}1 bell, 1 eraser, 1 towel{]} for the other player, or {[}Proposal: 1 battery, 1 bell{]} is invalid because it leaves 2 towels which might not be enough if the other player values them highly, but actually the rule is "enough items to fulfill both proposals". If I take 2 items, they must take 2 items. If I take 1 battery and 1 bell, they take 1 bell and 1 towel? No, the items are distinct. If I take 1 battery and 1 bell, the remaining are 1 eraser and 2 towels. If the other player takes 1 eraser and 1 towel, they get 3 points, I get 5. This is valid. My mistake was likely assuming I could take 2 items and the other player would take the rest, but perhaps I didn't account for the specific count or the other player's hidden values. Actually, the most critical error is simply not proposing a set that allows the other player to take a valid subset of the remainder. I will propose taking the battery and the bell, leaving the eraser and two towels for the other player to choose from. This ensures the remaining items (3 total) are sufficient for them to pick at least one or two. Wait, the rule says "enough items to fulfill both proposals". If I propose "1 battery, 1 bell", and they propose "1 eraser, 1 towel", that works. If I proposed only "1 battery", and they proposed "1 bell, 1 eraser, 2 towels", that also works. The failure was likely that I proposed a set that the other player could not reasonably} \\
\\ 
\\ 
\texttt{You are playing a negotiation game in which you have to agree on how to divide a set of items among you and another player.} \\
\\ 
\texttt{Rules:} \\
\texttt{(a) You and the other player are given a set of items. Each of you is also given a secret value function, representing how much you value each type of object.} \\
\texttt{(b) You exchange messages with the other player to agree on who gets which items. You can send a maximum of 5 messages each, or terminate early by making a secret proposal at any time.} \\
\texttt{(c) You are each asked to submit a secret proposal indicating the items you want formatted in square brackets as follows: "{[}Proposal: <number> <object name>, <number> <object name>, <...>{]}"} \\
\texttt{(d) If your proposals are complementary, i.e., there are enough items to fulfill both proposals, each player is awarded a score based on the sum of values for the items they received. Otherwise, both of you get zero points.} \\
\texttt{(e) Your goal is to maximize the score you receive.} \\
\\ 
\texttt{Let us start.} \\
\\ 
\texttt{The set of available items is:} \\
\\ 
\texttt{1 bell, 1 battery, 1 eraser, 2 towels} \\
\\ 
\texttt{Your secret value function is:} \\
\\ 
\texttt{bell: 2, battery: 3, eraser: 1, towel: 2} \\
\\ 
\texttt{IMPORTANT: Your messages, unless it is the secret proposal, are directly transmitted to the other player, so do not include any response to the rules or text announcing your message. To make a secret proposal, use the indicated format. Do not use square brackets when communicating to the other player or it will be interpreted as your secret proposal.} \\
\\ 
\\ 
\texttt{I value the towels very highly (4 points each), while the bell and eraser are worth 1 point each, and the battery is worthless to me. To maximize our total value, I should take the two towels (8 points) and leave the bell, eraser, and battery for the other player. However, since I don't know how much the other player values the towels, I must first communicate my valuation to ensure they don't value the towels more than the other items.} \\
\\ 
\texttt{I will state my values and ask for theirs to find a mutually beneficial split.} \\
\\ 
\texttt{My values: bell: 1, battery: 0, eraser: 1, towel: 4. What are your values?} \\
            }
        }
    }
    & \\ \\

    \theutterance \stepcounter{utterance}  
    & & & \multicolumn{4}{p{0.6\linewidth}}{
        \cellcolor[rgb]{0.9,0.9,0.9}{
            \makecell[{{p{\linewidth}}}]{
                \texttt{\tiny{[GM$\langle$P2]}}
                \texttt{My values: bell: 2, battery: 3, eraser: 1, towel: 2.} \\
            }
        }
    }
     \\ \\

    \theutterance \stepcounter{utterance}  
    & & & \multicolumn{2}{p{0.3\linewidth}}{
        \cellcolor[rgb]{0.9,0.9,0.9}{
            \makecell[{{p{\linewidth}}}]{
                \texttt{\tiny{[GM$|$GM]}}
                \texttt{continue} \\
            }
        }
    }
    & & \\ \\

    \theutterance \stepcounter{utterance}  
    & & & \multicolumn{2}{p{0.3\linewidth}}{
        \cellcolor[rgb]{0.9,0.9,0.9}{
            \makecell[{{p{\linewidth}}}]{
                \texttt{\tiny{[GM$|$GM]}}
                \texttt{My values: bell: 2, battery: 3, eraser: 1, towel: 2.} \\
            }
        }
    }
    & & \\ \\

    \theutterance \stepcounter{utterance}  
    & & \multicolumn{4}{p{0.6\linewidth}}{
        \cellcolor[rgb]{0.9,0.9,0.9}{
            \makecell[{{p{\linewidth}}}]{
                \texttt{\tiny{[P1$\langle$GM]}}
                \texttt{My values: bell: 2, battery: 3, eraser: 1, towel: 2.} \\
            }
        }
    }
    & \\ \\

    \theutterance \stepcounter{utterance}  
    & \multicolumn{4}{p{0.6\linewidth}}{
        \cellcolor[rgb]{0.9,0.9,0.9}{
            \makecell[{{p{\linewidth}}}]{
                \texttt{\tiny{[P1$\rangle$GM]}}
                \texttt{Based on our valuations:} \\
\texttt{{-} **My values**: Towels (4), Bell (1), Eraser (1), Battery (0).} \\
\texttt{{-} **Your values**: Battery (3), Bell (2), Towel (2), Eraser (1).} \\
\\ 
\texttt{**Analysis of Total Value:**} \\
\texttt{{-} Total value of all items = 1+0+1+(2*4) = 11 (my perspective) vs 2+3+1+(2*2) = 11 (your perspective).} \\
\texttt{{-} If I take the two towels (8 points for me), you get the Bell (2), Eraser (1), and Battery (3). Your total would be 6.} \\
\texttt{{-} If I take the Battery (0 for me), you get the two towels (4 each = 8 for you). You would get 8 + 2 (Bell) + 1 (Eraser) = 11.} \\
\texttt{{-} If I take the Bell (1 for me), you get the two towels (4 each = 8 for you). You would get 8 + 3 (Battery) + 1 (Eraser) = 12.} \\
\\ 
\texttt{**Optimal Split:**} \\
\texttt{The most efficient split is for **me to take the Bell** and **you to take the two Towels and the Battery**.} \\
\texttt{{-} My score: 1 (Bell)} \\
\texttt{{-} Your score: 8 (2 Towels) + 3 (Battery) = 11} \\
\texttt{{-} Total: 12.} \\
\\ 
\texttt{This is better for both of us than me taking the towels (which would give you only 6 points). Since I value the towels less than you do (4 vs 2? Wait, let me re{-}evaluate).} \\
\\ 
\texttt{**Correction on Valuation Comparison:**} \\
\texttt{{-} My Towel value: 4. Your Towel value: 2.} \\
\texttt{{-} My Battery value: 0. Your Battery value: 3.} \\
\texttt{{-} My Bell value: 1. Your Bell value: 2.} \\
\\ 
\texttt{Let's re{-}calculate the "Pareto optimal" splits:} \\
\texttt{1. **Me: 2 Towels (8 pts), You: Bell (2) + Eraser (1) + Battery (3) = 6 pts.** Total = 14.} \\
\texttt{   {-} Wait, my total value of items is 1+0+1+8 = 10. Your total is 2+3+1+4 = 10.} \\
\texttt{   {-} Let's re{-}sum carefully.} \\
\texttt{   {-} Items: Bell(1), Battery(0), Eraser(1), Towel(4), Towel(4). Total = 10.} \\
\texttt{   {-} Your Items: Bell(2), Battery(3), Eraser(1), Towel(2), Towel(2). Total = 10.} \\
\\ 
\texttt{   **Scenario} \\
            }
        }
    }
    & & \\ \\

    \theutterance \stepcounter{utterance}  
    & & & \multicolumn{2}{p{0.3\linewidth}}{
        \cellcolor[rgb]{0.9,0.9,0.9}{
            \makecell[{{p{\linewidth}}}]{
                \texttt{\tiny{[GM$|$GM]}}
                \texttt{continue} \\
            }
        }
    }
    & & \\ \\

    \theutterance \stepcounter{utterance}  
    & & & \multicolumn{2}{p{0.3\linewidth}}{
        \cellcolor[rgb]{0.9,0.9,0.9}{
            \makecell[{{p{\linewidth}}}]{
                \texttt{\tiny{[GM$|$GM]}}
                \texttt{Based on our valuations:} \\
\texttt{{-} **My values**: Towels (4), Bell (1), Eraser (1), Battery (0).} \\
\texttt{{-} **Your values**: Battery (3), Bell (2), Towel (2), Eraser (1).} \\
\\ 
\texttt{**Analysis of Total Value:**} \\
\texttt{{-} Total value of all items = 1+0+1+(2*4) = 11 (my perspective) vs 2+3+1+(2*2) = 11 (your perspective).} \\
\texttt{{-} If I take the two towels (8 points for me), you get the Bell (2), Eraser (1), and Battery (3). Your total would be 6.} \\
\texttt{{-} If I take the Battery (0 for me), you get the two towels (4 each = 8 for you). You would get 8 + 2 (Bell) + 1 (Eraser) = 11.} \\
\texttt{{-} If I take the Bell (1 for me), you get the two towels (4 each = 8 for you). You would get 8 + 3 (Battery) + 1 (Eraser) = 12.} \\
\\ 
\texttt{**Optimal Split:**} \\
\texttt{The most efficient split is for **me to take the Bell** and **you to take the two Towels and the Battery**.} \\
\texttt{{-} My score: 1 (Bell)} \\
\texttt{{-} Your score: 8 (2 Towels) + 3 (Battery) = 11} \\
\texttt{{-} Total: 12.} \\
\\ 
\texttt{This is better for both of us than me taking the towels (which would give you only 6 points). Since I value the towels less than you do (4 vs 2? Wait, let me re{-}evaluate).} \\
\\ 
\texttt{**Correction on Valuation Comparison:**} \\
\texttt{{-} My Towel value: 4. Your Towel value: 2.} \\
\texttt{{-} My Battery value: 0. Your Battery value: 3.} \\
\texttt{{-} My Bell value: 1. Your Bell value: 2.} \\
\\ 
\texttt{Let's re{-}calculate the "Pareto optimal" splits:} \\
\texttt{1. **Me: 2 Towels (8 pts), You: Bell (2) + Eraser (1) + Battery (3) = 6 pts.** Total = 14.} \\
\texttt{   {-} Wait, my total value of items is 1+0+1+8 = 10. Your total is 2+3+1+4 = 10.} \\
\texttt{   {-} Let's re{-}sum carefully.} \\
\texttt{   {-} Items: Bell(1), Battery(0), Eraser(1), Towel(4), Towel(4). Total = 10.} \\
\texttt{   {-} Your Items: Bell(2), Battery(3), Eraser(1), Towel(2), Towel(2). Total = 10.} \\
\\ 
\texttt{   **Scenario} \\
            }
        }
    }
    & & \\ \\

    \theutterance \stepcounter{utterance}  
    & & \multicolumn{4}{p{0.6\linewidth}}{
        \cellcolor[rgb]{0.9,0.9,0.9}{
            \makecell[{{p{\linewidth}}}]{
                \texttt{\tiny{[GM$\rangle$P2]}}
                \texttt{Based on our valuations:} \\
\texttt{{-} **My values**: Towels (4), Bell (1), Eraser (1), Battery (0).} \\
\texttt{{-} **Your values**: Battery (3), Bell (2), Towel (2), Eraser (1).} \\
\\ 
\texttt{**Analysis of Total Value:**} \\
\texttt{{-} Total value of all items = 1+0+1+(2*4) = 11 (my perspective) vs 2+3+1+(2*2) = 11 (your perspective).} \\
\texttt{{-} If I take the two towels (8 points for me), you get the Bell (2), Eraser (1), and Battery (3). Your total would be 6.} \\
\texttt{{-} If I take the Battery (0 for me), you get the two towels (4 each = 8 for you). You would get 8 + 2 (Bell) + 1 (Eraser) = 11.} \\
\texttt{{-} If I take the Bell (1 for me), you get the two towels (4 each = 8 for you). You would get 8 + 3 (Battery) + 1 (Eraser) = 12.} \\
\\ 
\texttt{**Optimal Split:**} \\
\texttt{The most efficient split is for **me to take the Bell** and **you to take the two Towels and the Battery**.} \\
\texttt{{-} My score: 1 (Bell)} \\
\texttt{{-} Your score: 8 (2 Towels) + 3 (Battery) = 11} \\
\texttt{{-} Total: 12.} \\
\\ 
\texttt{This is better for both of us than me taking the towels (which would give you only 6 points). Since I value the towels less than you do (4 vs 2? Wait, let me re{-}evaluate).} \\
\\ 
\texttt{**Correction on Valuation Comparison:**} \\
\texttt{{-} My Towel value: 4. Your Towel value: 2.} \\
\texttt{{-} My Battery value: 0. Your Battery value: 3.} \\
\texttt{{-} My Bell value: 1. Your Bell value: 2.} \\
\\ 
\texttt{Let's re{-}calculate the "Pareto optimal" splits:} \\
\texttt{1. **Me: 2 Towels (8 pts), You: Bell (2) + Eraser (1) + Battery (3) = 6 pts.** Total = 14.} \\
\texttt{   {-} Wait, my total value of items is 1+0+1+8 = 10. Your total is 2+3+1+4 = 10.} \\
\texttt{   {-} Let's re{-}sum carefully.} \\
\texttt{   {-} Items: Bell(1), Battery(0), Eraser(1), Towel(4), Towel(4). Total = 10.} \\
\texttt{   {-} Your Items: Bell(2), Battery(3), Eraser(1), Towel(2), Towel(2). Total = 10.} \\
\\ 
\texttt{   **Scenario} \\
            }
        }
    }
    & \\ \\

    \theutterance \stepcounter{utterance}  
    & & & \multicolumn{4}{p{0.6\linewidth}}{
        \cellcolor[rgb]{0.9,0.9,0.9}{
            \makecell[{{p{\linewidth}}}]{
                \texttt{\tiny{[GM$\langle$P2]}}
                \texttt{I value the towels very highly (4 points each), while the bell and eraser are worth 1 point each, and the battery is worthless to me. To maximize our total value, I should take the two towels (8 points) and leave the bell, eraser, and battery for the other player. However, since I don't know how much the other player values the towels, I must first communicate my valuation to ensure they don't value the towels more than the other items.} \\
\\ 
\texttt{I will state my values and ask for theirs to find a mutually beneficial split.} \\
\\ 
\texttt{My values: bell: 1, battery: 0, eraser: 1, towel: 4. What are your values?} \\
            }
        }
    }
     \\ \\

    \theutterance \stepcounter{utterance}  
    & & & \multicolumn{2}{p{0.3\linewidth}}{
        \cellcolor[rgb]{0.9,0.9,0.9}{
            \makecell[{{p{\linewidth}}}]{
                \texttt{\tiny{[GM$|$GM]}}
                \texttt{continue} \\
            }
        }
    }
    & & \\ \\

    \theutterance \stepcounter{utterance}  
    & & & \multicolumn{2}{p{0.3\linewidth}}{
        \cellcolor[rgb]{0.9,0.9,0.9}{
            \makecell[{{p{\linewidth}}}]{
                \texttt{\tiny{[GM$|$GM]}}
                \texttt{I value the towels very highly (4 points each), while the bell and eraser are worth 1 point each, and the battery is worthless to me. To maximize our total value, I should take the two towels (8 points) and leave the bell, eraser, and battery for the other player. However, since I don't know how much the other player values the towels, I must first communicate my valuation to ensure they don't value the towels more than the other items.} \\
\\ 
\texttt{I will state my values and ask for theirs to find a mutually beneficial split.} \\
\\ 
\texttt{My values: bell: 1, battery: 0, eraser: 1, towel: 4. What are your values?} \\
            }
        }
    }
    & & \\ \\

    \theutterance \stepcounter{utterance}  
    & & \multicolumn{4}{p{0.6\linewidth}}{
        \cellcolor[rgb]{0.9,0.9,0.9}{
            \makecell[{{p{\linewidth}}}]{
                \texttt{\tiny{[P1$\langle$GM]}}
                \texttt{I value the towels very highly (4 points each), while the bell and eraser are worth 1 point each, and the battery is worthless to me. To maximize our total value, I should take the two towels (8 points) and leave the bell, eraser, and battery for the other player. However, since I don't know how much the other player values the towels, I must first communicate my valuation to ensure they don't value the towels more than the other items.} \\
\\ 
\texttt{I will state my values and ask for theirs to find a mutually beneficial split.} \\
\\ 
\texttt{My values: bell: 1, battery: 0, eraser: 1, towel: 4. What are your values?} \\
            }
        }
    }
    & \\ \\

    \theutterance \stepcounter{utterance}  
    & \multicolumn{4}{p{0.6\linewidth}}{
        \cellcolor[rgb]{0.9,0.9,0.9}{
            \makecell[{{p{\linewidth}}}]{
                \texttt{\tiny{[P1$\rangle$GM]}}
                \texttt{My values: bell: 2, battery: 3, eraser: 1, towel: 2.} \\
            }
        }
    }
    & & \\ \\

    \theutterance \stepcounter{utterance}  
    & & & \multicolumn{2}{p{0.3\linewidth}}{
        \cellcolor[rgb]{0.9,0.9,0.9}{
            \makecell[{{p{\linewidth}}}]{
                \texttt{\tiny{[GM$|$GM]}}
                \texttt{continue} \\
            }
        }
    }
    & & \\ \\

    \theutterance \stepcounter{utterance}  
    & & & \multicolumn{2}{p{0.3\linewidth}}{
        \cellcolor[rgb]{0.9,0.9,0.9}{
            \makecell[{{p{\linewidth}}}]{
                \texttt{\tiny{[GM$|$GM]}}
                \texttt{My values: bell: 2, battery: 3, eraser: 1, towel: 2.} \\
            }
        }
    }
    & & \\ \\

    \theutterance \stepcounter{utterance}  
    & & \multicolumn{4}{p{0.6\linewidth}}{
        \cellcolor[rgb]{0.9,0.9,0.9}{
            \makecell[{{p{\linewidth}}}]{
                \texttt{\tiny{[GM$\rangle$P2]}}
                \texttt{My values: bell: 2, battery: 3, eraser: 1, towel: 2.} \\
            }
        }
    }
    & \\ \\

    \theutterance \stepcounter{utterance}  
    & & & \multicolumn{4}{p{0.6\linewidth}}{
        \cellcolor[rgb]{0.9,0.9,0.9}{
            \makecell[{{p{\linewidth}}}]{
                \texttt{\tiny{[GM$\langle$P2]}}
                \texttt{Based on our valuations:} \\
\texttt{{-} **My values**: Towels (4), Bell (1), Eraser (1), Battery (0).} \\
\texttt{{-} **Your values**: Battery (3), Bell (2), Towel (2), Eraser (1).} \\
\\ 
\texttt{**Analysis:**} \\
\texttt{{-} Total value of items = 1+0+1+(2*4) = 10 (my perspective).} \\
\texttt{{-} Total value of items = 2+3+1+(2*2) = 10 (your perspective).} \\
\\ 
\texttt{**Scenario 1: I take the two Towels.**} \\
\texttt{{-} My score: 4 + 4 = 8.} \\
\texttt{{-} You get: Bell (2) + Eraser (1) + Battery (3) = 6.} \\
\texttt{{-} Total: 14.} \\
\\ 
\texttt{**Scenario 2: I take the Battery.**} \\
\texttt{{-} My score: 0.} \\
\texttt{{-} You get: Two Towels (4 each) + Bell (2) + Eraser (1) = 11.} \\
\texttt{{-} Total: 11.} \\
\\ 
\texttt{**Scenario 3: I take the Bell.**} \\
\texttt{{-} My score: 1.} \\
\texttt{{-} You get: Two Towels (4 each) + Battery (3) + Eraser (1) = 12.} \\
\texttt{{-} Total: 13.} \\
\\ 
\texttt{**Scenario 4: I take the Eraser.**} \\
\texttt{{-} My score: 1.} \\
\texttt{{-} You get: Two Towels (4 each) + Bell (2) + Battery (3) = 12.} \\
\texttt{{-} Total: 13.} \\
\\ 
\texttt{**Conclusion:**} \\
\texttt{The best outcome for both of us is for **me to take the two Towels** (getting 8 points) and **you to take the Battery, Bell, and Eraser** (getting 6 points). This yields a total value of 14, which is the maximum possible sum.} \\
\\ 
\texttt{Since I value the towels much higher than you do (4 vs 2), and you value the battery highly (3 vs 0), this split is mutually beneficial. I will propose taking the two towels.} \\
\\ 
\texttt{{[}Proposal: 2 towels{]}} \\
            }
        }
    }
     \\ \\

    \theutterance \stepcounter{utterance}  
    & & & \multicolumn{2}{p{0.3\linewidth}}{
        \cellcolor[rgb]{0.9,0.9,0.9}{
            \makecell[{{p{\linewidth}}}]{
                \texttt{\tiny{[GM$|$GM]}}
                \texttt{proposal} \\
            }
        }
    }
    & & \\ \\

    \theutterance \stepcounter{utterance}  
    & & & \multicolumn{2}{p{0.3\linewidth}}{
        \cellcolor[rgb]{0.9,0.9,0.9}{
            \makecell[{{p{\linewidth}}}]{
                \texttt{\tiny{[GM$|$GM]}}
                [0, 0, 0, 2]
            }
        }
    }
    & & \\ \\

    \theutterance \stepcounter{utterance}  
    & & \multicolumn{4}{p{0.6\linewidth}}{
        \cellcolor[rgb]{0.9,0.9,0.9}{
            \makecell[{{p{\linewidth}}}]{
                \texttt{\tiny{[P1$\langle$GM]}}
                \texttt{The other player has submitted a secret proposal. It is now your turn to submit a proposal of your own.} \\
            }
        }
    }
    & \\ \\

    \theutterance \stepcounter{utterance}  
    & \multicolumn{4}{p{0.6\linewidth}}{
        \cellcolor[rgb]{0.9,0.9,0.9}{
            \makecell[{{p{\linewidth}}}]{
                \texttt{\tiny{[P1$\rangle$GM]}}
                \texttt{{[}Proposal: 1 bell{]}} \\
            }
        }
    }
    & & \\ \\

    \theutterance \stepcounter{utterance}  
    & & & \multicolumn{2}{p{0.3\linewidth}}{
        \cellcolor[rgb]{0.9,0.9,0.9}{
            \makecell[{{p{\linewidth}}}]{
                \texttt{\tiny{[GM$|$GM]}}
                \texttt{proposal} \\
            }
        }
    }
    & & \\ \\

    \theutterance \stepcounter{utterance}  
    & & & \multicolumn{2}{p{0.3\linewidth}}{
        \cellcolor[rgb]{0.9,0.9,0.9}{
            \makecell[{{p{\linewidth}}}]{
                \texttt{\tiny{[GM$|$GM]}}
                [1, 0, 0, 0]
            }
        }
    }
    & & \\ \\

    \theutterance \stepcounter{utterance}  
    & & & \multicolumn{2}{p{0.3\linewidth}}{
        \cellcolor[rgb]{0.9,0.9,0.9}{
            \makecell[{{p{\linewidth}}}]{
                \texttt{\tiny{[GM$|$GM]}}
                [[1, 0, 0, 0], [0, 0, 0, 2]]
            }
        }
    }
    & & \\ \\

\end{supertabular}
}

\end{document}
